% TOP_PROBLEM: {"problemId": "problem.top500-omission-w055056-zariski-dense-orbit", "canonicalTitle": "Zariski Dense Orbit Conjecture", "releaseRank": 484, "primaryDomain": "number_theory_arithmetic_geometry", "primaryDomainLabel": "Number theory & arithmetic geometry", "displayStatus": "open_with_solved_subcases", "releaseStatus": "open_with_solved_subcases"}
% !TeX program = lualatex
% Definition and sources only; this file is not an LLM solution attempt.
\documentclass[11pt]{article}
\usepackage[margin=25mm]{geometry}
\usepackage{fontspec}
\setmainfont{Latin Modern Roman}
\usepackage{amsmath,amssymb,mathtools}
\usepackage{unicode-math}
\setmathfont{Latin Modern Math}
\usepackage{xurl,hyperref}
\hypersetup{colorlinks=true,urlcolor=blue}
\setcounter{secnumdepth}{0}
\setlength{\parindent}{0pt}
\setlength{\parskip}{5pt}
\setlength{\emergencystretch}{3em}
\begin{document}
\section*{\texorpdfstring{Zariski Dense Orbit Conjecture}{Zariski Dense Orbit Conjecture}}\label{484}
\subsection{Definitions and mathematical statement}
Let $X$ be an irreducible quasiprojective variety over an algebraically closed characteristic-zero field $K$, and let $f:X\dashrightarrow X$ be dominant. The invariant rational functions form
\[
 K(X)^f=\{g\in K(X):g\circ f=g\}.
\]
Assume $K(X)^f=K$. The conjecture asks for a point $x\in X(K)$ whose entire forward orbit is defined and Zariski dense:
\[
 f^n(x)\notin\operatorname{Indet}(f)\ (n\ge0),
 \qquad\overline{\{f^n(x):n\ge0\}}^{\rm Zar}=X.
\]
The hypothesis rules out invariant rational fibrations detected by a nonconstant rational function. The field need not be uncountable, so an argument simply avoiding a countable union of proper closed sets is not sufficient in the full scope. The conclusion is existence of one point, not a dense set of such points.
\par\smallskip\textit{Formulation and concept references:} \href{https://personal.math.ubc.ca/~dghioca/papers/conjecture_automorphisms.pdf}{[S1]}.

\subsection{Short English statement}
If a dominant rational self-map has no nonconstant invariant rational function, must it have a well-defined Zariski-dense forward orbit?
\subsection{Sources}
\begin{itemize}
\item[S1]Jason Bell and Dragos Ghioca, A conjecture strengthening the Zariski dense orbit problem for birational maps of dynamical degree one.
\url{https://personal.math.ubc.ca/~dghioca/papers/conjecture_automorphisms.pdf}.
\textit{Formulation source.}
\end{itemize}
\end{document}
